\documentclass[10pt,border=2mm]{standalone}
\usepackage[T1]{fontenc}
\usepackage{figurestyle}
\input{primitives.tex}
\input{molecules.tex}
\input{numbers.tex}
\pgfplotsset{every axis/.append style={font=\fontsize{8}{10}\selectfont,line width=.4pt,tick style={line width=.4pt}}}
\begin{document}
\begin{tikzpicture}[plate]
\path[use as bounding box] (0,0) rectangle (15,33);
\node[sub,anchor=west] at (.2,32.6) {Figure 2, variant A: a comb of fibres over a slice of $\conf$; three configurations stand on the slice; bars are the two components};
% ================================================================ FIG 2 A
\begin{scope}[shift={(0,23.6)}]
  \begin{scope}[shift={(1.6,4.0)}]
    % the slice of C as a baseline (theta from 80 to 180 at delta = 0), x = (theta - 80)/100 * 10.4
    \draw[heavy] (-.2,0)--(10.6,0);
    \node[sub,anchor=west] at (10.7,0) {a slice of $\conf$: $\delta=0$, $\theta$ from $80^\circ$ to $180^\circ$};
    \foreach \th in {80,90,...,180} {\pgfmathsetmacro{\x}{(\th-80)/100*10.4}\draw[fine] (\x,0)--(\x,2.6);}
    \node[sub,anchor=south] at (5.2,2.65) {the same $\spin$ over every $X$};
    % bars: a (blue) and b (red) from the zero mark at height 1.3; scale .9 cm per unit
    \foreach \th in {80,90,...,180} {\pgfmathsetmacro{\x}{(\th-80)/100*10.4}
      \pgfmathsetmacro{\av}{exp(-((\th-104.5)/14.0)^2/2)}\pgfmathsetmacro{\bv}{0.9*((\th-118.0)/16.0)*exp(-((\th-118.0)/16.0)^2/2)}
      \draw[fine] (\x-.08,1.3)--(\x+.08,1.3);
      \draw[line width=2.2pt,draw=ColNitrogen] (\x-.07,1.3)--(\x-.07,{1.3+.9*\av});
      \draw[line width=2.2pt,draw=ColOxygen] (\x+.07,1.3)--(\x+.07,{1.3+.9*\bv});}
    \foreach \th/\k/\av/\bv in {\sampleThetaone/1/\sampleAone/\sampleBone,\sampleThetatwo/2/\sampleAtwo/\sampleBtwo,\sampleThetathree/3/\sampleAthree/\sampleBthree} {
      \pgfmathsetmacro{\x}{(\th-80)/100*10.4}
      \draw[fine] (\x,0)--(\x,2.6); \draw[fine] (\x-.08,1.3)--(\x+.08,1.3);
      \draw[line width=2.2pt,draw=ColNitrogen] (\x-.07,1.3)--(\x-.07,{1.3+.9*\av});
      \draw[line width=2.2pt,draw=ColOxygen] (\x+.07,1.3)--(\x+.07,{1.3+.9*\bv});
      \fill[PlateInk] (\x,0) circle[radius=1.6pt];
      \begin{scope}\molsolid\def\radH{.17}\def\radX{.27}\pic[scale=.75] at (\x,-.75) {mol3d-water-sample-\k};\end{scope}
      \node[lab] at (\x,-1.5) {$X_\k$};
      \node[sub] at (\x,-1.85) {$\Psi(X_\k)=\av\,\xi\,\bv\,\eta$};}
    \node[sub,anchor=west] at (10.7,1.3) {zero of $\spin$};
    \node[sub,anchor=west] at (10.7,.9) {\textcolor{ColNitrogen}{\rule{2.2pt}{6pt}} $a(X)$\quad \textcolor{ColOxygen}{\rule{2.2pt}{6pt}} $b(X)$};
  \end{scope}
  \node[sub,anchor=north west,text width=5.2cm] at (10.0,2.2) {$\Psi(X)=a(X)\,\xi+b(X)\,\eta$, with $\xi=\uparrow\downarrow\omega$, $\eta=\downarrow\uparrow\omega$ orthonormal in $\spin=\mathbb C^2\otimes\mathbb C^2\otimes\mathbb C^1$; one continuous representative, illustrative};
  \begin{axis}[at={(1.6cm,-.2cm)},width=8.6cm,height=3.6cm,xmin=80,xmax=180,ymin=-.75,ymax=1.3,ytick={-0.5,0,0.5,1},xtick={80,100,120,140,160,180},
    xlabel={$\theta$},xlabel style={at={(axis description cs:1,0)},anchor=west,xshift=3pt},axis lines=left,axis line style={-},tick align=outside,clip=false]
    \addplot[ColNitrogen,line width=.9pt] table[x=theta,y=a] {component-functions.dat};
    \addplot[ColOxygen,line width=.9pt,densely dashed] table[x=theta,y=b] {component-functions.dat};
    \addplot[PlateInk,line width=.6pt,dash pattern=on .6pt off 1.3pt] table[x=theta,y=n2] {component-functions.dat};
    \addplot[only marks,mark=*,mark size=1.6pt,PlateInk] coordinates {(\sampleThetaone,{\sampleAone*\sampleAone+\sampleBone*\sampleBone}) (\sampleThetatwo,{\sampleAtwo*\sampleAtwo+\sampleBtwo*\sampleBtwo}) (\sampleThetathree,{\sampleAthree*\sampleAthree+\sampleBthree*\sampleBthree})};
    \node[sub,anchor=west] at (axis cs:181,0.02) {$|a|^2+|b|^2$};
  \end{axis}
  \node[lab,anchor=north west] at (10.0,.6) {$\|\Psi\|^2=\displaystyle\int_{\conf}\bigl(|a|^2+|b|^2\bigr)\,dX<\infty$};
  \node[sub,anchor=north west,text width=4.8cm] at (10.0,-.6) {the dots are $\|\Psi(X_k)\|^2$ at the three configurations; $\Psi\sim\Phi$ when the integral of $\|\Psi-\Phi\|^2$ vanishes};
\end{scope}
% ================================================================ FIG 2 B
\node[sub,anchor=west] at (.2,19.2) {Figure 2, variant B: three fibres only, larger; the plot carries the continuum};
\begin{scope}[shift={(0,11.6)}]
  \begin{scope}[shift={(1.2,4.4)}]
    \draw[heavy] (-.2,0)--(8.4,0); \node[sub,anchor=west] at (8.5,0) {slice of $\conf$};
    \foreach \th/\k/\av/\bv/\x in {\sampleThetaone/1/\sampleAone/\sampleBone/1.0,\sampleThetatwo/2/\sampleAtwo/\sampleBtwo/4.1,\sampleThetathree/3/\sampleAthree/\sampleBthree/7.2} {
      \draw[fine] (\x,0)--(\x,2.9); \node[sub,anchor=south] at (\x,2.95) {$\spin$};
      \draw[fine] (\x-.1,1.4)--(\x+.1,1.4);
      \draw[line width=3pt,draw=ColNitrogen] (\x-.1,1.4)--(\x-.1,{1.4+1.2*\av});
      \draw[line width=3pt,draw=ColOxygen] (\x+.1,1.4)--(\x+.1,{1.4+1.2*\bv});
      \node[sub,anchor=west] at (\x+.22,{1.4+1.2*\av}) {$a=\av$}; \node[sub,anchor=west] at (\x+.22,{1.4+1.2*\bv}) {$b=\bv$};
      \fill[PlateInk] (\x,0) circle[radius=1.6pt];
      \begin{scope}\molsolid\pic[scale=.9] at (\x,-.9) {mol3d-water-sample-\k};\end{scope}
      \node[lab] at (\x,-1.8) {$X_\k$}; \node[sub] at (\x,-2.15) {$\theta=\th^\circ$};}
  \end{scope}
  \begin{axis}[at={(10.3cm,1.9cm)},width=4.6cm,height=3.4cm,xmin=80,xmax=180,ymin=-.75,ymax=1.3,ytick={-0.5,0,0.5,1},xtick={80,130,180},
    xlabel={$\theta$},xlabel style={at={(axis description cs:1,0)},anchor=west,xshift=3pt},axis lines=left,axis line style={-},tick align=outside,clip=false]
    \addplot[ColNitrogen,line width=.9pt] table[x=theta,y=a] {component-functions.dat};
    \addplot[ColOxygen,line width=.9pt,densely dashed] table[x=theta,y=b] {component-functions.dat};
    \addplot[PlateInk,line width=.6pt,dash pattern=on .6pt off 1.3pt] table[x=theta,y=n2] {component-functions.dat};
  \end{axis}
  \node[lab,anchor=north west] at (10.3,1.3) {$\|\Psi\|^2=\displaystyle\int_{\conf}\bigl(|a|^2+|b|^2\bigr)\,dX$};
\end{scope}
% ================================================================ FIG 3 A
\node[sub,anchor=west] at (.2,9.6) {Figure 3, variant A: spin slots in label order, leaders from the digits; slots swap with the labels};
\begin{scope}[shift={(0,5.2)}]
  \tikzset{slot/.style={draw=PlateInk,line width=.45pt,minimum width=.62cm,minimum height=.48cm,inner sep=1pt,fill=white,font=\fontsize{8}{9}\selectfont}}
  \foreach \x/\la/\lb/\pic/\sa/\sb/\cap in {2.0/1/2/mol3d-water-X/1/2/{$X$},6.6/2/1/mol3d-water-X/2/1/{$\sigma X=XP_\sigma$},11.2/1/2/mol3d-water-minusX/1/2/{$E^*X=-X$}} {
    \begin{scope}[shift={(\x,2.6)}]\molsolid\def\molA{\la}\def\molB{\lb}\pic[scale=1.0] at (0,0) {\pic};\end{scope}
    \node[lab] at (\x,1.35) {\cap};
    \node[slot] (s1) at (\x-.9,.45) {$\xi_{\sa}$}; \node at (\x-.45,.45) {$\otimes$};
    \node[slot] (s2) at (\x,.45) {$\xi_{\sb}$}; \node at (\x+.45,.45) {$\otimes$};
    \node[slot] (s3) at (\x+.9,.45) {$\xi_3$};
    \node[sub] at (\x,-.1) {slots $1,2,3$ in label order};}
  \draw[harrow] (3.6,2.6)--(5.0,2.6); \node[lab,anchor=south] at (4.3,2.7) {$\sigma=(12)$}; \node[sub,anchor=north] at (4.3,2.5) {relabel};
  \draw[harrow] (8.2,2.6)--(9.6,2.6); \node[lab,anchor=south] at (8.9,2.7) {$E^*$}; \node[sub,anchor=north] at (8.9,2.5) {invert};
  \node[lab,anchor=west] at (.4,-.8) {$\Psi(\sigma X)\sim\stat(\sigma)\,\sigma\act\Psi(X)$: the factors follow the labels, the sign $\stat(\sigma)=\prod_\alpha(\operatorname{sgn}\sigma_\alpha)^{2I_\alpha}$ sits outside};
  \node[lab,anchor=west] at (.4,-1.4) {$\Psi(E^*X)\sim\pm\Psi(X)$ on $\hilb^\pm$: parity is a separate character, $\stat$ ignores the star};
\end{scope}
% ================================================================ FIG 3 B
\node[sub,anchor=west] at (.2,2.9) {Figure 3, variant B: slots under the nuclei; in $\sigma X$ the row reads $\xi_2\,\xi_1\,\xi_3$; the inversion triad after Harter (1993) fig. 5.1.1};
\begin{scope}[shift={(0,-1.6)}]
  \tikzset{slot/.style={draw=PlateInk,line width=.45pt,minimum width=.6cm,minimum height=.46cm,inner sep=1pt,fill=white,font=\fontsize{8}{9}\selectfont}}
  \foreach \x/\la/\lb/\pic/\cap in {2.0/1/2/mol3d-water-X/{$X$},6.0/2/1/mol3d-water-X/{$\sigma X$},10.0/1/2/mol3d-water-minusX/{$E^*X$}} {
    \begin{scope}[shift={(\x,2.9)}]\molsolid\def\molA{\la}\def\molB{\lb}\pic[scale=1.0] at (0,0) {\pic};\end{scope}
    \node[lab] at (\x,1.5) {\cap};}
  % slots under the nuclei of X (H1 left, H2 right, O below centre), of sigma X (left is 2), of E*X (inverted: H's below)
  \foreach \x/\l/\r in {2.0/1/2,6.0/2/1} {\node[slot] at (\x-.85,.8) {$\xi_{\l}$}; \node[slot] at (\x+.85,.8) {$\xi_{\r}$}; \node[slot] at (\x,.3) {$\xi_3$};}
  \node[slot] at (10.0-.85,.3) {$\xi_1$}; \node[slot] at (10.0+.85,.3) {$\xi_2$}; \node[slot] at (10.0,.8) {$\xi_3$};
  \node[sub] at (2.0,-.15) {$\Psi(X)=\xi_1\otimes\xi_2\otimes\xi_3$}; \node[sub] at (6.0,-.15) {$\Psi(\sigma X)\sim\stat(\sigma)\,\xi_2\otimes\xi_1\otimes\xi_3$}; \node[sub] at (10.0,-.15) {$\Psi(E^*X)\sim\pm\Psi(X)$};
  \draw[harrow] (3.5,2.9)--(4.5,2.9); \node[lab,anchor=south] at (4.0,3.0) {$\sigma=(12)$};
  \draw[harrow] (7.5,2.9)--(8.5,2.9); \node[lab,anchor=south] at (8.0,3.0) {$E^*$};
  % inversion triad on a sphere
  \begin{scope}[shift={(13.2,2.6)}]
    \hatom{0}{0}{.75}{PlateInk!6}
    \draw[harrow,line width=.6pt] (0,0)--(0,1.05); \node[sub,anchor=south] at (0,1.07) {$z$};
    \draw[harrow,line width=.6pt] (0,0)--(1.05,0); \node[sub,anchor=west] at (1.07,0) {$y$};
    \draw[harrow,line width=.6pt] (0,0)--(-.5,-.6); \node[sub,anchor=north east] at (-.5,-.6) {$x$};
    \draw[hidden,-{Stealth[length=1.6mm]}] (0,0)--(0,-1.0); \draw[hidden,-{Stealth[length=1.6mm]}] (0,0)--(-1.0,0); \draw[hidden,-{Stealth[length=1.6mm]}] (0,0)--(.48,.58);
    \node[sub,anchor=north,text width=3.0cm,align=center] at (0,-1.25) {$E^*$ sends the triad to its negative (dashed): right-handed becomes left-handed, which no rotation does};
  \end{scope}
\end{scope}
\end{tikzpicture}
\end{document}
